% Chapter 7 can be included from a larger thesis with \input or \include.
% Conclusions follow Chapter 6's checksum-verified 60 Hz v5 formal cohort.
% Historical DSP-only, smoke and nominal-clock reports are not pooled with it.

\chapter{Conclusions and Future Work}
\label{chap:conclusions_future_work}

This dissertation investigated how streaming music can guide the selection,
timing and transition of humanoid dance drawn from a finite motion library.
The resulting Humanoid Matcher connects offline AIST++ preparation and G1
retargeting with online music retrieval, stateful motion selection and
continuous kinematic playback. The experiments establish a working,
inspectable integration with measurable rhythm alignment and a validated 60 Hz
operating point, while also identifying an unmet 120 Hz target. This chapter summarises the
work, states the limits of those conclusions and identifies the developments
needed for broader musical coverage and physical robot execution.

\section{Summary of the Work}
\label{sec:ch7_summary}

\paragraph{Data preparation and motion representation.}
The offline pipeline established a traceable association between music,
human choreography and target-robot trajectories. Audio windowing, robust
amplitude normalisation and source-motion analysis produced comparable music
descriptors and velocity-valley key poses. SMPL sequences were converted through
GMR into named-joint, floating-base trajectories for the 29-DoF Unitree G1.
Versioned artefacts, source and model hashes, validated caching and MuJoCo
preflight made catalogue admission explicit. The recorded build contains
411 admitted motions linked to 60 music identities and 348 audio segments.
These checks establish representation and sampled kinematic validity for the
stored clips; they do not certify every subsequent retiming or blend.

\paragraph{Hierarchical music retrieval.}
The final catalogue uses the Discogs EffNet ONNX representation, comprising a
1,280-dimensional embedding and 400 style activations, together with a
24-dimensional DSP rhythm--timbre descriptor. The retargeting artefacts use
pipeline version 4. The DSP-only embedding remains an alternative backend,
requiring a separately built catalogue, rather than the representation evaluated
as the final system. Online queries use six seconds of accumulated audio and
the same extraction and normalisation contract as the reference catalogue.
Segment evidence is pooled at track level, filtered by genre and abstention
rules, and passed to motion ranking. Strict leave-one-music-out evaluation in
Chapter~6 obtained genre Recall@1 of 0.317 and Recall@5 of 0.783. Thus the
representation supplies useful shortlist evidence, although its leading result
is frequently incorrect for unseen music.

\paragraph{Motion compatibility and stable selection.}
The system separates musical similarity from suitability of the associated
movement. Tempo feasibility, key-pose density and motion activity refine the
track-based ranking before selection. Consecutive-winner evidence and a score
margin provide relevance hysteresis, while a bounded-hold policy permits
variation among similarly relevant motions. Prepared-candidate requirements
separate selection intent from execution readiness. These mechanisms make
repetition, rejection and switching decisions observable. Their implementation
is a contribution of the work, but the replay ablations do not establish a
statistically supported benefit for every component after multiplicity
correction; in particular, a diversity benefit cannot be inferred merely from
the presence of a rotation policy.

\paragraph{Beat-aware timing and gradual phase control.}
The timing controller distinguishes detected pulses from events accepted for
motion alignment. Relative beat contrast and authored key-pose spacing guide
beat acceptance, and bounded speed adaptation consumes phase error gradually
instead of resetting the sampler instantaneously. In the matched causal short
suite, harmonic beat alignment score (BAS) was 0.306 under authored timing and
0.313 with the full controller. The source-averaged gain was 0.0067, with a
95\% bootstrap interval of [-0.0078, 0.0228] and a Holm-adjusted $p$-value of
1.0000. The formal experiment therefore does not establish a timing benefit
under the specified delivered-beat metric. Nor does it establish accuracy
against independently annotated musical beats or an improvement in perceived
dance quality.

\paragraph{Continuous switching and output conditioning.}
The transition mechanism searches neighbourhoods of salient target poses using
joint position, velocity, foot-contact, root-state and salience costs. Root
alignment, smoothstep blending, quaternion interpolation and atomic state
transfer connect successive moving trajectories. One-shot completion avoids
an implicit last-to-first seam, while joint limiting and configurable collision
handling condition the output. These mechanisms provide an explicit route from
a selected motion to a displayed pose. The final-output audit verifies that all
135 formal causal runs satisfy the registered joint-range, speed, acceleration
and clearance constraints at 60 Hz. This remains a kinematic contract rather
than evidence of dynamic balance or hardware safety.

\paragraph{Real-time implementation and evaluation.}
Retrieval and candidate preparation execute in bounded background workers,
while the foreground loop consumes completed results and advances playback.
The evaluation records rankings, beat delivery, preparation, transitions,
final poses and actual output clocks. It also separates older preanalysed
file replay from the causal supplement, which contains 130 formal short runs
and five long runs. All 135 pass the registered real-time, readiness and final-
output gates at 60 Hz. A separate 120 Hz pilot passes only 6/21 runs, so the
original frequency is not a reliable operating point on this platform. The
overall achievement is therefore a reproducible retrieval-and-playback system
with an empirically validated 60 Hz configuration and a measured capacity limit.

\section{Limitations}
\label{sec:ch7_limitations}

\paragraph{Kinematic playback and physical feasibility.}
The current player writes root pose and joint positions directly to MuJoCo
state and evaluates forward kinematics. It does not apply a torque-controlled
tracking policy or solve closed-loop balance. Grounding and binary contact
features describe geometric support, without enforcing contact forces, friction,
actuator capacity or disturbance recovery. Consequently, a displayed G1
trajectory may remain dynamically infeasible even when its joint positions and
sampled geometry are valid. No physical Unitree G1 deployment has been evaluated
in this dissertation.

\paragraph{Operating-point and final-output limits.}
The formal 60 Hz cohort passes all registered output constraints: maximum speed
is 15.404 rad/s, maximum acceleration is 1673.415 rad/s$^2$, and both residual
clearance and final joint-range violation counts are zero. Ready-pool underruns
and hold-last events are also zero. These results apply to audited kinematic
poses on the recorded machine. They do not establish torque feasibility,
balance, contact stability or physical-robot safety.

For the causal full short suite, the median run-level control-work p99 is
4.657 ms; the five long runs have a median of 8.724 ms. All 135 formal runs
remain below the 16.67 ms p99 limit and 0.01 deadline-miss-ratio limit, although
the worst long run reaches 14.483 ms and 0.00693 respectively. In contrast, the
separate 120 Hz pilot passes only 6/21 runs: its worst p99 is 15.572 ms against
an 8.33 ms period and its worst miss ratio is 0.18544. The 60 Hz evidence is
therefore empirical soft-real-time evidence for this platform, not a hard-
real-time guarantee or evidence that the original 120 Hz target has been met.

\paragraph{Finite coverage and unequal genre performance.}
The matcher adapts existing clips and cannot supply a movement vocabulary
absent from the catalogue. Its coverage depends on the available AIST++ pairs,
their retargeting quality and manually designed compatibility weights. Musical
similarity is also an indirect proxy for choreographic suitability. Genre
Recall@1 ranged from 1.000 for house to 0.000 for ballet jazz, street jazz and
krump in the six-identities-per-genre evaluation. Broad Discogs-to-AIST++ label
mappings and early genre restriction can propagate a classification error into
the motion candidate set. These results reveal substantial class differences,
while the small number of identities per class limits generalisation.

\paragraph{Metre assumptions and weak-beat behaviour.}
The switch clock groups every four accepted alignment beats into a logical
bar. It does not estimate musical metre or identify downbeats. Since beat
acceptance depends on motion feasibility, rejected subdivisions can also change
the relationship between this clock and the underlying musical pulse. The rule
is therefore poorly suited to triple metre, including waltz, and cannot promise
phrase-aligned switching even for all quadruple-metre music. The observed waltz
retrieval errors should not, however, be attributed to this scheduling rule
without an experiment that isolates metre from retrieval.

For ambient or otherwise weakly rhythmic input, persistent negative-style and
beat-quality gates provide conservative abstention. The ambient CC0 recording
was rejected in 96\% of evaluated windows, but one recording is insufficient
evidence of robust handling of the category. Holding the active motion or
returning towards authored timing can remain musically inappropriate when the
audio becomes quiet or loses its pulse. In particular, reducing secondary pose
modulation does not necessarily stop the underlying authored trajectory. The
system lacks a fully specified transition from dancing to low-activity behaviour
and back, calibrated across diverse weak-beat recordings.

\paragraph{Engineering proxies and experimental coverage.}
No user study assesses naturalness, expressiveness or preference. BAS measures
alignment to delivered detector events, while contact proxies, entropy,
coverage and smoothness each capture only part of dance quality. Reconstructed
SMPL FID/Div additionally depend on trace interpolation, source-motion reuse
and cyclic reference construction. Near-zero reference variance makes some
standardised geometric results ill-conditioned, so they cannot support an
aesthetic ranking. The ten external CC0 recordings and synthetic negatives
offer limited acoustic diversity; headless file-input trials do not reproduce
microphone reverberation, clipping or room noise. Repeated seeds improve
repeatability evidence without creating independent musical sources, and the
causal long runs repeat only one stitched source. The causal supplement also
retests only full versus authored timing, so the complete replay ablation matrix
cannot be treated as causal online validation.

\section{Future Work}
\label{sec:ch7_future_work}

\paragraph{Extend the validated operating envelope.}
The immediate systems priority is to preserve the coherent final-output contract
while reducing foreground work enough to evaluate 120 Hz again. A short-horizon
constrained trajectory update could incorporate collision correction within the
same temporal problem as joint limiting. Its feasibility and fallback behaviour
should continue to be evaluated from actual output timestamps rather than from
an intermediate limiter report. Profiling should guide reductions in repeated
grounding work and especially foreground collision cost, with prepared data
reused under an explicit model-and-version contract. New comparisons should
freeze the machine, backend,
window settings and audit configuration, and counterbalance run order to reduce
background-load confounding. Explicit query-slot identifiers and a persistent
candidate identifier across match, pending, readiness and blend events would
also make end-to-end response attribution more reliable.

\paragraph{Metre and downbeat-aware control.}
A causal beat, downbeat and metre estimator could replace the fixed four-event
clock with an estimate of musical position. The scheduler should distinguish
the underlying beat sequence from the subset used for motion alignment, so
skipping an infeasible accent does not redefine the bar. Uncertain metre should
lead to a documented beat-level fallback rather than an assumed downbeat.
Evaluation should include triple metre, compound metre, tempo changes and weak
percussion, using independently annotated beat and downbeat times alongside
detector delivery latency. A complementary phase study should cover at least
three distinct motions per genre, stratified by tempo and key-pose density;
multiple windows or seeds of one motion should not count as additional motions.

\paragraph{Learned transition-entry scoring.}
The current weighted entry cost offers a useful baseline for learning which
candidate state produces a successful transition. Training examples could
combine current and target pose, velocity, contact, root motion and musical
context with measured transition outcomes and human pairwise preferences.
A learned ranker could then capture interactions that a fixed weighted sum
misses, while explicit feasibility checks continue to determine admissibility.
Evaluation should hold out source motions and music identities and compare the
ranker with the existing cost on both perceived continuity and final-output
constraints. An improved prediction score alone would not establish a better
executed transition.

\paragraph{Motion-fragment retrieval and composition.}
The existing system can enter partway through a selected clip, but retrieval
and activity metadata remain primarily clip-level. A fragment catalogue could
instead index phrases or contact-consistent movement units, each with local
music descriptors, entry and exit states, duration and feasible tempo range.
Searching compatible fragment sequences would permit finer responses to changes
within a song and reuse a larger combination of existing motifs. A sparse
transition graph or bounded-horizon search could balance music relevance,
contact compatibility and motif continuity. Such an extension needs new checks
for the joins and composed sequence: validating each fragment separately would
not establish that their concatenation is feasible. Diversity should remain
subordinate to coherent phrasing rather than encourage rapid switching.

\paragraph{Contact-constrained tracking and balance.}
Foot-contact states should become constraints on execution rather than only
terms in entry ranking. A subsequent controller could combine support-phase
information, foot placement and floating-base state estimation with whole-body
inverse dynamics, model-predictive control or a learned tracking policy. Its
role would be to track the selected reference subject to contact and actuator
constraints, reducing motion intensity or delaying a transition when necessary.
Dynamics-based simulation should evaluate tracking error, foot slip, support
transitions, torque saturation and recovery from disturbances. This would test
a different and stronger claim than the present kinematic preview.

\paragraph{Physical Unitree G1 deployment.}
Hardware integration should follow validated closed-loop simulation and a
measured final-output contract. The music matcher would supply motion references
to the tracking controller rather than directly assign the physical root state.
Initial trials could progress from supported, low-amplitude movements to
freestanding sequences and then music-driven transitions. They should record
microphone acquisition, communication and actuator-response delays separately,
as well as tracking failures and interventions. This staged evaluation would
reveal how acoustic conditions, state-estimation error and actuator limits
change the behaviour observed in headless file-input experiments.

\paragraph{User preference and weak-beat adaptation.}
A preference-aware selector could allow users to favour particular styles,
energy levels or degrees of repetition within the feasible candidate set.
Pairwise judgements of musical fit, transition naturalness and overall appeal
would support both preference learning and a direct evaluation of dance quality.
Weak-beat behaviour should similarly become a selectable, explicit response:
for example, a controlled change to a low-activity motion followed by a gradual
return when rhythmic evidence recovers. Additional real ambient, speech, noise
and mixed recordings are needed to calibrate rejection and recovery separately.
User preferences should influence the choice among feasible motions while
preserving execution constraints, and subjective gains should be reported
alongside relevance, continuity and timing measurements.

The work provides a basis for these extensions by making the chain from musical
evidence to robot pose traceable. Its current results support finite-library
retrieval and a 60 Hz kinematic operating point, but do not establish that causal
phase control improves BAS or that 120 Hz is reliable. They demonstrate why
musical responsiveness, output feasibility and real-time compliance must be
evaluated separately. Resolving their interaction is the next step towards a
humanoid that can perform continuously and expressively under physical control.
